Existing dominant methods for audio generation include Generative Adversarial Networks (GANs) and diffusion-based methods like Flow Matching. GANs suffer from slow convergence
during training, while diffusion methods require multi-step inference that introduces considerable computational overhead. In this work, we introduce Flow2GAN, a two-stage framework that combines Flow Matching training for learning generative capabilities with GAN fine-tuning for efficient few-step inference. Specifically, given audio's unique properties, we first improve Flow Matching for audio modeling through: 1) reformulating the objective as endpoint estimation, avoiding velocity estimation difficulties when involving empty regions; 2) applying spectral energy-based loss scaling to emphasize perceptually salient quieter regions. 
Building on these Flow Matching adaptations, we demonstrate that a further stage of lightweight GAN fine-tuning enables us to obtain few-step (e.g., 1/2/4 steps) generators that produce high-quality audio. In addition, we develop a multi-branch network architecture that processes Fourier coefficients at different time-frequency resolutions, which improves the modeling capabilities compared to prior single-resolution designs. 
Experimental results indicate that our Flow2GAN delivers high-fidelity audio generation from Mel-spectrograms or discrete audio tokens, achieving highly favorable quality-efficiency trade-offs compared to existing state-of-the-art GAN-based and Flow Matching-based methods.
Online demo samples are available at \url{https://flow2gan.github.io},
and the source code is released at \url{https://github.com/k2-fsa/Flow2GAN}. 
% Experimental results indicate that our Flow2GAN delivers high-fidelity audio generation from Mel-spectrograms or discrete audio tokens, outperforming existing state-of-the-art GAN-based and Flow Matching-based methods in both quality and efficiency.

% Deep generative models have demonstrated remarkable progress in audio generation, enabling reconstructing high-quality waveforms from compressed representations. Although Generative Adversarial Networks (GANs) have emerged as a dominant approach with one-step inference, they often suffer from slow convergence and potential mode collapse during training. Recent diffusion-based methods like Flow Matching, which offer stable training and powerful generative capabilities, can achieve GAN-competitive performance, but their multi-step inference introduces considerable computational overhead. In this work, we introduce Flow2GAN, a novel framework that combines the strengths of both paradigms, featuring two-stage training and multi-resolution network design, where the first stage employs Flow Matching for learning robust generative capabilities, and the second stage applies GAN fine-tuning for finer-grained generation with efficient one-/two-step inference. Specifically, given audio's unique properties, we first improve Flow Matching for audio modeling through: 1) reformulating the objective as endpoint estimation, avoiding velocity estimation difficulty in empty regions; 2) applying spectral energy-based loss scaling to emphasize perceptually salient quieter regions. Building on these Flow Matching adaptations, we demonstrate that fast GAN fine-tuning enables us to obtain one-step and two-step generators that produce high-quality audio. In addition, we develop a multi-branch network that processes Fourier coefficients at different time-frequency resolutions, showing substantial improvements in modeling capabilities over prior single-resolution designs. Extensive experiments indicate that our Flow2GAN framework delivers high-fidelity audio generation from Mel-spectrograms or discrete acoustic tokens, outperforming existing state-of-the-art GAN-based and Flow Matching-based methods in both quality and efficiency. 


% Deep generative models have demonstrated remarkable progress in audio generation, enabling reconstructing high-quality waveforms from compressed representations. Although Generative Adversarial Networks (GANs) have emerged as a dominant approach with one-step inference, utilizing sophisticated discriminators to capture multi-grained audio details, they often suffer from slow convergence and potential mode collapse during training. Recent diffusion-based methods like Flow Matching, which offer stable training and powerful generative capabilities, can achieve GAN-competitive performance when enhanced with multi-band waveform processing techniques, but their multi-step inference introduces considerable computational overhead. In this work, we introduce Flow2GAN, a two-stage training strategy that combines the strengths of both paradigms, where Flow Matching training in the first stage provides robust generative capabilities, and GAN fine-tuning in the second stage guides the model toward finer-grained generation with efficient one-/two-step inference.  Specifically, considering the unique characteristics of audio signals, we first improve Flow Matching for audio modeling through: 1) reformulating the objective as endpoint estimation, avoiding velocity estimation difficulty in empty regions; 2) applying spectral energy-based loss scaling to emphasize perceptually salient quieter regions. Building on these Flow Matching adaptations, we demonstrate that lightweight GAN fine-tuning enables us to obtain one-step and two-step generators that produce high-quality audio. Additionally, we develop a multi-branch network architecture that processes Fourier spectral coefficients at different time-frequency resolutions, showing substantial improvements in modeling capabilities over prior single-branch design. Extensive experiments show that our Flow2GAN framework delivers high-fidelity audio generation, outperforming existing state-of-the-art GAN-based and Flow Matching-based methods in both quality and efficiency.

